Fix \(P\in\mathfrak P_{\mathcal D}^{\mathrm{obs}}\) and \(B>0\).  By Definition~\(\mathrm{def{:}observable\text{-}bridge\text{-}law\text{-}class}\), Assumption~\(\mathrm{ass{:}nondegenerate\text{-}components}\), and positivity of the unit costs, every \(V_{E,d},V_{O,d},c_{E,d},c_{O,d}\) is finite and strictly positive.  Theorem~\(\mathrm{thm{:}budget\text{-}profile}\) therefore gives
\[
 0<v_d
 =\left(\sqrt{c_{E,d}V_{E,d}}+
          \sqrt{c_{O,d}V_{O,d}}\right)^2
 <\infty
 \qquad(d\in\mathcal D).                                      \tag{10}
\]
Because \(\mathcal D\) is finite and nonempty, \(\mathcal O_P\) is finite and nonempty: it contains \((d,d)\) for every \(d\).  Thus the maximum in (1) exists, is finite, and is at least one, proving
\[
 1\le C_{\mathrm{ord}}(P)<\infty.                             \tag{11}
\]

Suppose first that the implication on the left of (2) holds with a constant \(C\).  For every \((d,d')\in\mathcal O_P\), it gives \(v_d\le Cv_{d'}\); division by the positive denominator in (10) yields \(v_d/v_{d'}\le C\).  Maximizing over \(\mathcal O_P\) gives \(C_{\mathrm{ord}}(P)\le C\).  Conversely, if \(C\ge C_{\mathrm{ord}}(P)\) and \(p_d\le p_{d'}\), then \((d,d')\in\mathcal O_P\), so
\[
 v_d\le C_{\mathrm{ord}}(P)v_{d'}\le Cv_{d'}.
\]
This proves (2).

Finiteness of \(\mathcal D\) also guarantees that the diagnostic minimizer set and all minima in (3) are attained.  Choose
\[
 d_{\mathrm p}^*\in
 \arg\min_{d\in\mathcal D_{\mathrm{pred}}(P)}v_d,
 \qquad
 d^*\in\arg\min_{d\in\mathcal D}v_d.
\]
Since \(d_{\mathrm p}^*\) minimizes \(p_d\) over \(\mathcal D\), one has \(p_{d_{\mathrm p}^*}\le p_{d^*}\), hence \((d_{\mathrm p}^*,d^*)\in\mathcal O_P\).  It follows from (10)--(11) that
\[
 1\le
 \frac{v_{d_{\mathrm p}^*}}{v_{d^*}}
 \le C_{\mathrm{ord}}(P).                                    \tag{12}
\]
By Definition~\(\mathrm{def{:}predictive\text{-}benchmark}\), the middle term in (12) is simultaneously \(C_{\mathrm{sel}}(P)\) and \(\rho_{\mathrm{pred}}(P)\), proving (3).

Let \(\varnothing\ne\mathcal Q\subseteq\mathfrak P_{\mathcal D}^{\mathrm{obs}}\).  If a finite constant \(C\) works for the best-tie rule at every \(P\in\mathcal Q\), exactness of (3) gives \(\rho_{\mathrm{pred}}(P)\le C\) for every such \(P\), so its supremum is finite.  Conversely, if that supremum is finite, the supremum itself is a common valid constant.  Replacing (3) by the equivalence (2) proves the identical assertion for all pairwise prediction-order implications and \(\sup_{P\in\mathcal Q}C_{\mathrm{ord}}(P)\).  This also proves that no causal condition defining the smaller class \(\mathfrak P_{\mathcal D}\) is needed for the exact lawwise comparisons.

Now take a numerical realization of \(\mathsf H_{\mathrm{comp}}\) satisfying (4)--(7).  Positivity of \(Q_{R,d}\) on \(\mathcal X_d(P)\) makes every ratio in (6) well-defined.  From the definitions of the infimum and supremum, for every \(d\in\mathcal D\) and \(x\in\mathcal X_d(P)\),
\[
 \Lambda_-(P)Q_{R,d}(x)
 \le F_d(x)
 \le\Lambda_+(P)Q_{R,d}(x).                                \tag{13}
\]
Moreover, these constants are sharp on the chosen cones.  Indeed, any common lower-envelope constant \(L\) satisfying \(LQ_{R,d}(x)\le F_d(x)\) for all \(d,x\) must obey
\[
 L\le\min_{d\in\mathcal D}\inf_{x\in\mathcal X_d(P)}
       \frac{F_d(x)}{Q_{R,d}(x)}=\Lambda_-(P),
\]
and any common upper-envelope constant \(U\) must analogously satisfy \(U\ge\Lambda_+(P)\).  Consequently a strictly positive finite two-sided envelope exists exactly when (7) holds, and (13) is the weakest such envelope in the stated sharp-constant sense.

For any \((d,d')\in\mathcal O_P\), equations (4) and (13) yield
\[
\begin{aligned}
 v_d
 &\le\Lambda_+(P)Q_{R,d}\{x_d(P)\}
   =\Lambda_+(P)Bp_d\\
 &\le\Lambda_+(P)Bp_{d'}
   =\Lambda_+(P)Q_{R,d'}\{x_{d'}(P)\}\\
 &\le\frac{\Lambda_+(P)}{\Lambda_-(P)}
       F_{d'}\{x_{d'}(P)\}
   =\frac{\Lambda_+(P)}{\Lambda_-(P)}v_{d'}.
\end{aligned}                                                \tag{14}
\]
The division is legitimate by \(\Lambda_-(P)>0\).  Maximizing (14) over \(\mathcal O_P\), and then applying (3), proves (8).

It remains to justify (9) and the generalized-eigenvalue qualification.  Put
\[
 u=c_{E,d}Q_{E,d}(x),
 \qquad
 w=c_{O,d}Q_{O,d}(x),
\]
so \(u,w\ge0\).  When \(u,w>0\), direct algebra gives, for every \(0<\alpha<1\),
\[
 \frac{u}{\alpha}+\frac{w}{1-\alpha}
 -(\sqrt u+\sqrt w)^2
 =\frac{\{(1-\alpha)\sqrt u-\alpha\sqrt w\}^2}
        {\alpha(1-\alpha)}\ge0.                              \tag{15}
\]
Equality holds at \(\alpha=\sqrt u/(\sqrt u+\sqrt w)\).  If exactly one of \(u,w\) vanishes, the same infimum is obtained by letting \(\alpha\) approach the appropriate endpoint; if both vanish, both sides are zero.  Hence (15) proves (9) in all cases.

The profiling operation need not preserve quadraticity.  For example, on \(\mathbb R^2\), with unit costs, \(Q_E(x)=x_1^2\) and \(Q_O(x)=x_2^2\), formula (5) becomes
\[
 F(x)=(|x_1|+|x_2|)^2.
\]
It violates the parallelogram identity because, for the coordinate vectors \(e_1,e_2\),
\[
 F(e_1+e_2)+F(e_1-e_2)=8
 \ne4=2F(e_1)+2F(e_2).
\]
Thus a profiled numerator is not in general a quadratic form and therefore is not in general represented by a single generalized Rayleigh quotient.  At a fixed \(\alpha\), the numerator in (9) is quadratic, so a largest generalized eigenvalue can furnish an upper full-space envelope when \(Q_{R,d}\) is positive definite; it does not furnish the positive lower envelope required in (13).  Finally, equality in (8) would require equality throughout (14), hence common-space comparability, normalization of the lower envelope, and law-realized directions that simultaneously saturate the relevant envelopes and diagnostic order.  These are additional conditions, and neither observable operator bijectivity nor the diagnostic ordering supplies them.  This completes the proof.